% Part: first-order-logic
% Chapter: semantics
% Section: intro-semantics

\documentclass[../../../include/open-logic-section]{subfiles}

\begin{document}

\olfileid{fol}{syn}{its}

\olsection{Inleiding}

De regels waarmee uitdrukkingen betekenis krijgen, noemen we
semantiek. Centraal staat de vraag wanneer een formule waar is in
!!a{structure}. !!^a{structure} legt vast wat de bouwstenen van de taal
betekenen. !!a{domain} is een niet-lege verzameling objecten. De
kwantoren gaan over de objecten in dit domein. Aan !!{constant}s worden
elementen uit het domein gekoppeld. De !!{function}s worden gekoppeld
aan functies die objecten uit het !!{domain} weer aan objecten uit
datzelfde domein koppelen. De !!{predicate}s worden gekoppeld aan
relaties op de objecten van het !!{domain}. Samen vormen het
!!{domain} en al deze koppelingen aan de basiswoordenschat
!!a{structure}. Ook !!^{variable}s kunnen in !!{formula}s staan. Om de
betekenis vast te leggen, koppelen we daarom ook aan die variabelen
!!{element}s van het !!{domain}. Dat heet een toekenning aan
variabelen. De relatie ``waar in/onder'' brengt al deze gegevens bij
elkaar. !!^a{formula} kan waar zijn in !!a{structure}~$\Struct{M}$ bij
!!a{variable} toekenning~$s$. We schrijven dat als
$\Sat{M}{!A}[s]$. Die relatie leggen we met inductie stap voor stap
vast volgens de opbouw van~$!A$. Bij $(!A \land !B)$ gebruiken we
bijvoorbeeld de tabellen met waarheidswaarden: of de samengestelde
formule waar is, hangt af van de vraag of $!A$ en ~$!B$ wel of niet
waar zijn. De !!{variable} toekenning maakt niet uit als de
!!{formula}~$!A$ !!a{sentence} is, dus zonder vrije variabelen. We
kunnen dan gewoon zeggen dat !!{sentence}s in !!{structure}s waar zijn
of niet.

Met de relatie $\Sat{M}{!A}$ voor !!{sentence}s kunnen we drie
basisbegrippen precies definiëren: geldigheid, semantisch gevolg en
vervulbaarheid. !!^a{sentence} is geldig, $\Entails !A$, als zij in
iedere !!{structure} waar is. We schrijven voor een verzameling
!!{sentence}s $\Gamma \Entails !A$ als iedere !!{structure} die alle
!!{sentence}s in~$\Gamma$ waar maakt, ook~$!A$ waar maakt. Dan volgt de
laatste zin semantisch uit die verzameling. Een verzameling
!!{sentence}s is vervulbaar als minstens één !!{structure} alle
!!{sentence}s erin tegelijk waar maakt. !!{formula}s zijn stap voor
stap opgebouwd. De regels die bepalen wanneer ze waar zijn, volgen op
hun beurt de opbouw van !!{formula}s. Daarom kunnen we met inductie
eigenschappen van de semantiek bewijzen en verbanden tussen deze
begrippen aantonen.
\end{document}
